%% kapa: exam=1 round=23 year=2012 period=1 form=A subject=economics subject_ko=경제원론 seq=2 number=42
%% source: Q-Net 감정평가사 기출문제 (cst00309) artlSeq=5052578 file=raw/1008/5052578/2012년 제23회 감정평가사 제1차시험 기출문제(1교시 A형).pdf
%% conversion: kordoc 4.24.51 -> tools/gen1.py (auto); review=proofread-sonnet; flags=none
\begin{question}{42}
제품의 가격이 10원이고, 노동 한 단위의 가격은 5원, 자본 한 단위의 가격은 15원이다. 기업 A의 노동의 한계생산이 3이고, 자본의 한계생산은 1일 때, 현재 생산수준에서 비용극소화를 위한 방법으로 옳은 것은? (단, 모든 시장은 완전경쟁시장이고, 노동과 자본의 한계생산은 체감한다.)
\begin{choices}
  \item 노동의 투입량은 늘리고, 자본의 투입량은 줄일 것이다.
  \item 노동의 투입량은 줄이고, 자본의 투입량은 늘릴 것이다.
  \item 노동과 자본 모두 투입량을 늘릴 것이다.
  \item 노동과 자본 모두 투입량을 줄일 것이다.
  \item 노동과 자본의 투입량을 그대로 유지할 것이다.
\end{choices}
\end{question}
\answer{1}
